% Bloco unico, como a Elsevier exige, e ate 250 palavras. Cada "fix" declarado
% com o status que a inferencia no nivel do pais sustenta (parecer de painel,
% rodada 2, P0).
\begin{abstract}

Um gestor ambiental municipal pergunta a um modelo de linguagem qual é o padrão
vinculante de PM\textsubscript{2,5} em seu país, e nada na resposta diz se ela
corresponde ao registro oficial. Propomos um protocolo para essa conferência e
relatamos o que ele encontra. As respostas são pontuadas contra registros
oficiais, por código sempre que a resposta correta é um valor publicado e por um
painel de juízes de três fornecedores nos demais casos, em 25 países, 14
configurações de acesso e quatro idiomas (\ncellspt{} células pontuadas).
Nenhuma das três mitigações que um órgão consegue implantar sem construir nada
funciona, e a mais intuitiva prejudica. Perguntar no idioma do próprio país
reduz a acurácia em \nnativa{} pontos percentuais (modelo misto com intercepto
aleatório por país, $p=\nhtwomixedcountryp$; negativa em \nhtwocneg{} de nove
países) e torna o modelo \nverifdiffabs{} pontos menos propenso a devolver
qualquer valor conferível contra o registro. Declarar o papel oficial equivale a
não declarar nada, em até dois pontos percentuais. O ajuste regional de
instrução em 7B não melhora o modelo-base de que partiu. A acurácia é mais baixa
na tarefa com consequência administrativa, recordar o valor vinculante
(\ntasktoneshort{} contra \nfloorrestshort{} para síntese e recomendação), e um
terço das respostas erradas cita um valor da própria escada oficial do país. Os
países do Sul Global devolvem o valor do registro \ncodegapabs{} pontos menos
vezes que os do Norte com o país como unidade; as estimativas pontuais são
maiores onde a resposta é um fato nacional, mas todo intervalo por tarefa
inclui~1.

\end{abstract}
